% Part: sets-functions-relations
% Chapter: sets
% Section: subsets

\documentclass[../../../include/open-logic-section]{subfiles}

\begin{document}

\olfileid{sfr}{set}{sub}
\olsection{উপসেট ও ঘাত সেট}

\begin{explain}
প্রায়ই আমাদের সেটগুলির তুলনা করতে হবে। এমন তুলনার একটি সহজ ধরন হল:
\emph{একটি সেটে যা কিছু আছে, অন্য সেটেও তা আছে}। এই পরিস্থিতি এতটাই
গুরুত্বপূর্ণ যে, এর জন্য আমরা কিছু নতুন প্রতীক প্রবর্তন করব।
\end{explain}

\begin{defn}[উপসেট]
কোনও সেট $A$-এর প্রতিটি !!{element} যদি $B$-এরও !!a{element} হয়,
তবে $A$-কে $B$-এর একটি \emph{উপসেট} বলি এবং $A \subseteq B$ লিখি।
$A$ যদি $B$-এর উপসেট না হয়, তবে $A \not\subseteq B$ লিখি।
$A \subseteq B$ কিন্তু $A \neq B$ হলে আমরা $A \subsetneq B$ লিখি
এবং বলি, $A$ হল $B$-এর একটি \emph{প্রকৃত উপসেট}।
\end{defn}

\begin{ex}
প্রতিটি সেট নিজেরই একটি উপসেট, আর $\emptyset$ প্রতিটি সেটের উপসেট।
জোড় স্বাভাবিক সংখ্যাগুলির সেট স্বাভাবিক সংখ্যাগুলির সেটের উপসেট।
এ ছাড়া, $\{ a, b \} \subseteq \{ a, b, c \}$। কিন্তু
$\{ a, b, e \}$ সেটটি $\{ a, b, c \}$-এর উপসেট নয়।
\end{ex}

\begin{ex}
$2$ সংখ্যাটি পূর্ণসংখ্যাগুলির সেটের !!{element}; অন্যদিকে, জোড় সংখ্যাগুলির
সেট পূর্ণসংখ্যাগুলির সেটের উপসেট। তবে কোনও সেট অন্য একটি সেটের
!!a{element} এবং উপসেট, \emph{দুই-ই} হতে পারে। যেমন,
$\{0\} \in \{0, \{0\}\}$ এবং $\{0\} \subseteq \{0, \{0\}\}$।
\end{ex}

সদস্যভিত্তিক সমতার নীতি সেটের অভিন্নতার একটি শর্ত দেয়: $A = B$ হবে
যদি এবং কেবল যদি $A$-এর প্রতিটি !!{element} $B$-এরও !!a{element} হয়
এবং বিপরীত দিক থেকেও একই কথা সত্য হয়। ``উপসেট''-এর সংজ্ঞা অনুযায়ী
$A \subseteq B$ ঠিক এই শর্তের প্রথম অর্ধাংশটিই বোঝায়:
$A$-এর প্রতিটি !!{element} $B$-এরও !!a{element}। অবশ্য $A$ ও $B$-এর
স্থান বদলালেও সংজ্ঞাটি প্রযোজ্য: অর্থাৎ, $B \subseteq A$ হবে যদি এবং
কেবল যদি $B$-এর প্রতিটি !!{element} $A$-এরও !!a{element} হয়।
এটিই সদস্যভিত্তিক সমতার নীতির ``বিপরীত দিক থেকেও'' অংশটির অর্থ।
অন্য কথায়, এই নীতি থেকে পাই: সেট দুটি সমান হবে যদি এবং কেবল যদি
তারা পরস্পরের উপসেট হয়।

\begin{prop}
$A = B$ হবে যদি এবং কেবল যদি $A \subseteq B$ এবং $B \subseteq A$, উভয়ই হয়।
\end{prop}

এখন আরও কিছু উপযোগী প্রতীক প্রবর্তনের সুযোগ এসেছে। $A$ কখন $B$-এর
উপসেট হয়, তার সংজ্ঞায় আমরা বলেছিলাম ``$A$-এর প্রতিটি !!{element}
\dots'', এবং ``$\dots$''-এর জায়গায় বসিয়েছিলাম ``$B$-এর
!!a{element}''। কিন্তু এই \emph{ধরনের} অভিব্যক্তি এত ঘন ঘন ব্যবহৃত হয়
যে, এর জন্য কিছু আনুষ্ঠানিক প্রতীক প্রবর্তন করা সুবিধাজনক হবে।

\begin{defn}\ollabel{forallxina}
$(\forall x \in A)\phi$ হল $\forall x(x \in A \lif
\phi)$-এর সংক্ষিপ্ত রূপ। একইভাবে, $(\exists x \in A)\phi$ হল
$\exists x(x \in A \land \phi)$-এর সংক্ষিপ্ত রূপ।
\end{defn}

এই প্রতীক ব্যবহার করে বলতে পারি, $A \subseteq B$ হবে যদি এবং কেবল যদি
$(\forall x \in A)x \in B$ হয়।

এবার এক বিশেষ ধরনের সেট বিবেচনা করব: একটি প্রদত্ত সেটের সমস্ত উপসেটের সেট।

\begin{defn}[ঘাত সেট]
কোনও সেট~$A$-এর সমস্ত উপসেট নিয়ে গঠিত সেটকে
\emph{$A$-এর ঘাত সেট} বলে এবং $\Pow{A}$ লেখা হয়।
  \[
    \Pow{A} = \Setabs{B}{B \subseteq A}
  \]
\end{defn}

\begin{ex}
$\{ a, b, c \}$-এর সমস্ত সম্ভাব্য উপসেট কী কী? সেগুলি হল:
$\emptyset$, $\{a \}$, $\{b\}$, $\{c\}$, $\{a, b\}$, $\{a, c\}$,
$\{b, c\}$, $\{a, b, c\}$। এই সমস্ত উপসেটের সেট হল
$\Pow{\{a,b,c\}}$:
\[
\Pow{\{ a, b, c \}} = \{\emptyset, \{a \}, \{b\}, \{c\}, \{a, b\},
\{b, c\}, \{a, c\}, \{a, b, c\}\}
\]
\end{ex}

\begin{prob}
$\{a, b, c, d\}$-এর সমস্ত উপসেটের তালিকা তৈরি করো।
\end{prob}

\begin{prob}
দেখাও যে, $A$-এর $n$টি !!{element}s থাকলে $\Pow{A}$-এর
$2^n$টি !!{element}s থাকে।
\end{prob}

\end{document}
